% language=us signal=signal

% \nologbuffering

\usemodule[present-boring,abbreviations-logos]

\definelayer
  [extratext]
  [preset=rightbottom,
   height=\textheight,
   width=\textwidth]

\setupbackgrounds
  [text]
  [background=extratext]

\setuptyping
  [align=hangright,
   keeptogether=paragraph,
   numbering=]

\startdocument
  [title={balancing mvls},
   banner={from one to the other},
   location={bachotex\enspace {\bf 2025}\enspace meeting}]

\starttitle[title=Disclaimer]

This is work in progress. We hope to be stable around the \CONTEXT\ meeting so
what we give here is an impression. The primitives shown are meant to be used in
mechanisms that bring them together.

In this talk we go:

\startitemize[packed]
\startitem from paragraphs \stopitem
\startitem via pages \stopitem
\startitem and columns \stopitem
\startitem to mvls \stopitem
\startitem as well as balancing \stopitem
\startitem and discuss what is involved \stopitem
\stopitemize

What follow is mostly an outline for myself.

\stoptitle

\starttitle[title=The context]

\startitemize

\startitem
    In \LUAMETATEX\ we have an extended, multi-pass par builder, the mechanism
    that breaks lines into paragraphs.
\stopitem

\startitem
    We also extended the page builder a bit, the mechanism that collects and
    split of content.
\stopitem

\startitem
    Alongside these we added a new mechanism that relates these two in terms of
    solution space.
\stopitem

\stopitemize

\stoptitle

\starttitle[title={Paragraphs}]

In the par builder the engine deals with:

\startitemize[columns,two,packed]
    \startitem orphan penalties       : lone words on the last line \stopitem
    \startitem toddler penalties      : single glyphs at the margin \stopitem
    \startitem twin penalties         : similar words at the margin \stopitem
    \startitem spacing glue           : glue between words \stopitem
    \startitem punctuation glue       : glue around punctuation \stopitem
    \startitem l r m boxes            : repeated content at breaks \stopitem
    \startitem various math           : all kind of constraints \stopitem
    \startitem (user) leaders         : adaptive content \stopitem
    \startitem glyph expansion        : stretching and shrinking \stopitem
    \startitem character protrusion   : escaping into margins \stopitem
    \startitem discretionaries        : end and begin line snippets \stopitem
    \startitem hyhenation (penalties) : breaks within words \stopitem
    \startitem optional content       : multi-pass wiggle room \stopitem
    \startitem shapes & indentation   : determines width \stopitem
\stopitemize

\startitemize[columns,two,packed]
    \startitem widow penalties        : between last lines  \stopitem
    \startitem club penalties         : between first lines \stopitem
    \startitem shape penalties        : between shape lines \stopitem
    \startitem broken penalties       : between hyphenated lines \stopitem
    \startitem interline penalties    : between lines \stopitem
\stopitemize

\startitemize[columns,two,packed]
    \startitem spacing (stretch & shrink) : criteria driven  \stopitem
    \startitem emergency stretch          : just in case \stopitem
    \startitem left & right margins       : determines width \stopitem
    \startitem left & right hanging       : determines width \stopitem
    \startitem initial left & right       : determines width \stopitem
    \startitem final left & right         : determines width \stopitem
\stopitemize

\stoptitle

\starttitle[title={Pages}]

In the page builder the engine deals with:

\startitemize[columns,two,packed]
    \startitem vsize             : the main constraint \stopitem
    \startitem inserts           : need to be bound \stopitem
    \startitem glue              : can ease constraints \stopitem
    \startitem penalties         : drives the routine \stopitem
    \startitem adjusted content  : used for special purposes \stopitem
    \startitem marks             : synchronized content \stopitem
    \startitem vz wiggle room    : less constraints \stopitem
    \startitem top & bottom skip : for aligning the bottom \stopitem
    \startitem extra goal        : optional wiggle room \stopitem
    \startitem feedback loop     : permits column construction \stopitem
\stopitemize

\stoptitle

\starttitle[title={Columns}]

See Mikaels presentation. We have different solutions for different situations:

\definehighlight[kindofbold][style=\bf\glyphweight100]

\startitemize
\startitem
    \kindofbold {Simple columns:} Collect content in a vbox and let \TEX\ do the
    splitting.
\stopitem
\startitem
    \kindofbold {Mixed columns:} Set the vsize to n times the text height and
    when we have a full page split the collected content. This is sensitive for
    spacing and decisions around page and column breaks. The splitting is mostly
    delegated to \LUA. ({\bf Boxed columns} are a variant of this.)
\stopitem
\startitem
    \kindofbold {Page columns:} Set the vsize to the text height and collect
    pages as usual but delay flushing till the number of columns is reached. The
    splitting is mostly done in \TEX\ using the regular builder.
\stopitem
\startitem
    \kindofbold {Column sets (old):} Change the vsize every time a pseudo page is
    split off. In \MKIV\ we let \LUA\ do some of the work.
\stopitem

\blank

\startitem
    \kindofbold {Column sets (new):} Collect content and define page template and
    at some point flush the content to fit those templates using multi-slot
    optimization.
\stopitem
\stopitemize

\stoptitle

\starttitle[title={Main vertical lists}]

\startitemize
\startitem
    In vertical mode the engine is either assembling a vbox (or equivalent like an insert)
    or adding to the main vertical list.
\stopitem
\startitem
    Splitting of pages from that vertical list happen as we go.
\stopitem
\startitem
    The vsplit routine is not the same as the page split routine.
\stopitem
\startitem
    In \CONTEXT\ we migrate deeply hidden material (inserts, marks, adjusts) to
    the outer level (in \LUAMETATEX\ natively or \MKIV\ via callbacks).
\stopitem
\startitem
    In \LUAMETATEX\ we can (efficiently) collect to different vertical lists
    without packaging (vboxing) or splitting (page break).
\stopitem
\stopitemize

\stoptitle

\starttitle[title={Balancing}]

\startitemize
\startitem
    This evolved out of a previous (\MKIV/\MKXL) columnset implementation.
\stopitem
\startitem
    We started with a similar approach as the par builder: choosing the best solution
    out of (possibly) many.
\stopitem

\startitem
    We started from the same routine but with specific features dropped. For a
    while we kept discretionary code.
\stopitem

\startitem
    Interfacing wise something like discretionaries makes no sense: there is no
    intuitive way to use something equivalent (conditional content at page
    breaks) also because decisions are made as one goes before breaking kicks in.
    So it was dropped (less code due to less active|/|passive breakpoints.
\stopitem

\startitem
    What did make sense, was exploring page shapes analogue to par shapes:
    heights (noflines), topskip and bottomskip conditions kick in. Shape entries
    relate to what we call slots in columnsets.
\stopitem

\startitem
    That set us on the path of trial runs and adaptive shapes. However this is
    done at the \TEX\ end in loops that can test conditions and results and
    retry.
\stopitem

\startitem
    At first we didn't want to do inserts but then realized that it was doable given some
    constraints: they are slot based.
\stopitem

\startitem
    Because we want to work on the grid, we need to take care of half-line spacing and
    synchronize at breaks. This introduces top and bottom optional content, basically
    dumb discretionaries but done differently.
\stopitem

\startitem
    In the end we also realized that trial runs made it possible to implement high
    performance balancing.
\stopitem

\stopitemize

\stoptitle

\starttitle[title=Activities: intercept]

We have to intercept a \quote {page stream}:

\starttyping
\beginmvl 1
    various content
\endmvl

\setbox\scratchboxone\flushmvl 1

\beginmvl
    index 2
  % prevdepth 4pt
    options \numexpr "1 + "4\relax
\relax
    various content
\endmvl
\stoptyping

\setlayerframed
    [extratext]
    [frame=on,offset=2mm,strut=no,align=normal]
    \bgroup
        \dorecurse{15}{\blackrule[height=.8sh,depth=.8sd,width=3cm,color=white]\par}
    \egroup

% \starttabulate[|Tr|||]
% \NC 0x\tohexadecimal\ignoreprevdepthmvloptioncode \NC ignore prevdepth \NC \type {\ignoreprevdepthmvloptioncode} \NC \NR
% \NC 0x\tohexadecimal\noprevdepthmvloptioncode     \NC no prevdepth     \NC \type {\noprevdepthmvloptioncode    } \NC \NR
% \NC 0x\tohexadecimal\discardtopmvloptioncode      \NC discard top      \NC \type {\discardtopmvloptioncode     } \NC \NR
% \NC 0x\tohexadecimal\discardbottommvloptioncode   \NC discard bottom   \NC \type {\discardbottommvloptioncode  } \NC \NR
% \stoptabulate

% \tracingmvl <number>
% \mvlcurrentlyactive

\stoptitle

\starttitle[title=Activities: shapes]

\starttyping
\balanceshape 3
    vsize      3\lineheight
    topskip    \strutht
    bottomskip \strutdp
next
    vsize      5\lineheight
    topskip    \strutht
    bottomskip \strutdp
next
    vsize      8\lineheight
    topskip    \strutht
    bottomskip \strutdp
\relax
\stoptyping

\starttyping
\setbox\scratchboxtwo\vbalance\scratchboxone
\stoptyping

\setlayerframed
    [extratext]
    [frame=on,offset=2mm,strut=no,align=normal]
    \bgroup % 15
        \hbox \bgroup
            \vbox \bgroup
                \framed[align=normal]\bgroup
                    \dorecurse{3}{\blackrule[height=.8sh,depth=.8sd,width=3cm,color=white]\par}
                \egroup\par
                \framed[align=normal]\bgroup
                    \dorecurse{5}{\blackrule[height=.8sh,depth=.8sd,width=3cm,color=white]\par}
                \egroup\par
                \framed[align=normal]\bgroup
                    \dorecurse{7}{\blackrule[height=.8sh,depth=.8sd,width=3cm,color=white]\par}
                    \dorecurse{1}{\blackrule[height=.8sh,depth=.8sd,width=3cm,color=extracolor]\par}
                \egroup\par
            \egroup
            \hskip2mm
            \vbox \bgroup
                \framed[align=normal]\bgroup
                    \dorecurse{3}{\blackrule[height=.8sh,depth=.8sd,width=3cm,color=white]\par}
                \egroup\par
                \framed[align=normal]\bgroup
                    \dorecurse{4}{\blackrule[height=.8sh,depth=.8sd,width=3cm,color=white]\par}
                    \dorecurse{1}{\blackrule[height=.8sh,depth=.8sd,width=3cm,color=extracolor]\par}
                \egroup\par
                \framed[align=normal]\bgroup
                    \dorecurse{8}{\blackrule[height=.8sh,depth=.8sd,width=3cm,color=white]\par}
                \egroup\par
            \egroup
        \egroup
    \egroup

\stoptitle

\starttitle[title=Activities: flush]

\starttyping
\hbox \bgroup
    \localcontrolledendless {%
        \ifvoid\scratchboxtwo
            \expandafter\quitloop
        \else
            \setbox\scratchbox\ruledhbox\bgroup
                \vbalancedbox\scratchboxtwo
            \egroup
            \vbox to 12\lineheight \bgroup
                \box\scratchbox
                \vfill
            \egroup
            \hskip1em
        \fi
    }\unskip
\egroup
\stoptyping

\stoptitle

\starttitle[title=Activities: trials]

\starttyping
\setbox\scratchboxtwo\vbalance\scratchboxone
    trial
\stoptyping

\setlayerframed
    [extratext]
    [frame=on,offset=2mm,strut=no,align=normal]
    \bgroup % 15
        \hbox \bgroup
            \vbox \bgroup
                \framed[align=normal]\bgroup
                    \dorecurse{3}{\strut\kern3cm\par}%
                \egroup\par
                \framed[align=normal]\bgroup
                    \dorecurse{5}{\strut\kern3cm\par}%
                \egroup\par
                \framed[align=normal]\bgroup
                    \dorecurse{7}{\strut\kern3cm\par}%
                \egroup\par
            \egroup
            \hskip2mm
            \vbox \bgroup
                \framed[align=normal]\bgroup
                    \dorecurse{3}{\strut\kern3cm\par}%
                \egroup\par
                \framed[align=normal]\bgroup
                    \dorecurse{4}{\strut\kern3cm\par}%
                \egroup\par
                \framed[align=normal]\bgroup
                    \dorecurse{8}{\strut\kern3cm\par}%
                \egroup\par
            \egroup
        \egroup
    \egroup

\stoptitle

\starttitle[title=Activities: breaks]

\starttyping
\balanceboundary 3 1\relax
\vskip\zeropoint
\balanceboundary 3 0\relax
\vskip\zeropoint
\balanceboundary 3 0\relax
\stoptyping

Using these with a proper callback we can go to a next slot. column, page or
spread.

% \starttabulate[|c|c|||]
% \BC first \BC second \BC action                                    \BC user interface          \NC \NR
% \NC 1     \NC 1 or 0 \NC goto next spread (1 initial, 0 follow up) \NC \type {\page[spread]}   \NC \NR
% \NC 2     \NC 1 or 0 \NC goto next page (idem)                     \NC \type {\page}           \NC \NR
% \NC 3     \NC 1 or 0 \NC goto next column (idem)                   \NC \type {\column}         \NC \NR
% \NC 4     \NC 1 or 0 \NC goto next slot (idem)                     \NC \type {\column[slot]}   \NC \NR
% \NC 5     \NC n      \NC next slot when more than n lines          \NC \type {\testroom[5]}    \NC \NR
% \NC 6     \NC s      \NC next slot when more than s scaled points  \NC \type {\testroom[80pt]} \NC \NR
% \stoptabulate

\stoptitle

\starttitle[title=Activities: marks]

Marks are kind of hidden in the resulting (split of) segments of the
balanced list. They can be set using \LUA\ calls:

\starttyping
node.direct.updatetopmarks()
node.direct.updatefirstmarks()
\stoptyping

and

\starttyping
node.direct.updatemarks(list)
node.direct.updatefirstandbotmarks(list)
\stoptyping

The pattern is:

\startitemize[packed]
\startitem set the top marks \stopitem
\startitem loop over columns and set marks \stopitem
\startitem set the first marks \stopitem
\stopitemize

\stoptitle

\starttitle[title=Activities: inserts]

\startitemize[packed]
\startitem We can assume a reasonable amount of notes. \stopitem
\startitem These are normally small with no (vertical) whitespace. \stopitem
\startitem Notes taking multiple lines may split. \stopitem
\startitem But we need to obey widow and club penalties. \stopitem
\startitem There can be math formulas but mostly inline. \stopitem
\startitem We need to keep them close to where they are referred from. \stopitem
\blank
\startitem We can ignore complex conflicting demands. \stopitem
\startitem As long as we get some result, we're fine. \stopitem
\startitem So users have to check what comes out. \stopitem
\startitem We don't assume fully automated unattended usage. \stopitem
\stopitemize

One can check for inserts with

\starttyping
<state> = \boxinserts <box>
\stoptyping

and fetch them with

\starttyping
<box>   = \vbalancedinsert <box> <class>
\stoptyping

possible states are: has_inserts, has_inserts_with_content, has_inserts_with_height

\stoptitle

\starttitle[title=Activities: discarding]

We need to deal with (for instance) half line spacing on the grid which is a bit
hackish but sort of generic. Kind of like this:

Trigger:

\starttyping
\par \vskip0pt \hrule discardable height 1sh depth 1sd width 1em \par
\stoptyping

Before:

\starttyping
\vskip0pt \vbox discardable {\hpack{\strut BEFORE}} \par
\stoptyping

After:

\starttyping
\vskip0pt \vbox discardable {\hpack{\strut AFTER}} \penalty -1 \par
\stoptyping

Handling:

\starttyping
\vbalanceddiscard 2       \relax % becomes a \nohrule
\vbalanceddiscard 2 remove\relax % disappears
\stoptyping

This mechanism is still evolving.

\stoptitle

\starttitle[title=Activities: penalties]

It is possible to discourage breaks at the end of a to be balanced list, like:

\starttyping
\balancefinalpenalties 6
    10000 9000 8000 7000 6000 5000
\relax
\stoptyping

\stoptitle

\starttitle[title=Activities: passes]

One can optimize the balancing with multiple passes, but we need to experiment a
bit more with this. We just mention variables to play with:

\startitemize[columns,three,packed]
\startitem threshold \stopitem
\startitem tolerance \stopitem
\startitem looseness \stopitem
\startitem adjdemerits \stopitem
\startitem originalstretch \stopitem
\startitem emergencystretch \stopitem
\startitem emergencyfactor \stopitem
\startitem emergencypercentage \stopitem
\stopitemize

This model is simpler than the one used in the par builder but we can relatively
easy extend it.

\stoptitle

\starttitle[title=Some examples]

We show some of the examples that we use for testing and figuring out additional
features. We really start from usability, not from some abstract cq.\ utopian
dream.

\stoptitle

\stopdocument
